% Unit 093 terjemahan Bahasa Indonesia dari chapter7.tex baris 1130--1279.
% Sumber: Methods of Algebra, Volume 2, commit
% 9a5803ff77dd3257484cb177f851a73770a59dd3 (CC BY 4.0).
% stable-unit-id: o014.aljabr2.chapter7.coalgebra-to-hopf-b
% Cakupan: sisa Bagian 7.5, dari sesudah bukti struktur monoidal kategori
% modul bialjabar hingga tepat sebelum Bagian 7.6 pada baris 1280.

% segment-id: o014.aljabr2.chapter7.coalgebra-to-hopf-b.p001
Untuk bialjabar $A$, kita masih dapat menanyakan apakah kategori monoidal
$A\dcate{Mod}$ memiliki struktur anyaman dan apakah struktur itu simetris.
Gagasan awal yang wajar ialah menggunakan struktur anyaman $c(M_1,M_2)$ dari
$\mathcal{V}$, tetapi pada umumnya ini bukan morfisme dalam $A\dcate{Mod}$.
Bahkan untuk kasus khusus $\mathcal{V}=\Bbbk\dcate{Mod}$, jawaban atas
pertanyaan ini sudah melibatkan konsep \emph{bialjabar kuasitriangular}, yang
tidak dapat dijelaskan secara singkat. Berikut hanya diberikan suatu contoh
yang begitu sederhana sehingga agak menyesatkan.

% segment-id: o014.aljabr2.chapter7.coalgebra-to-hopf-b.pr001
\begin{proposition}\label{prop:bialgebra-mod-symmetry}
	Misalkan $\mathcal{V}$ kategori monoidal simetris.
% segment-id: o014.aljabr2.chapter7.coalgebra-to-hopf-b.l001
	\begin{enumerate}[(i)]
% segment-id: o014.aljabr2.chapter7.coalgebra-to-hopf-b.i001
		\item Jika $A$ bialjabar kokomutatif di atas $\mathcal{V}$, maka kategori
		monoidal $A\dcate{Mod}$ dan $\cated{Mod}A$ menjadi kategori monoidal
		simetris terhadap
		$c(M_1,M_2):M_1\otimes M_2\rightiso M_2\otimes M_1$.

% segment-id: o014.aljabr2.chapter7.coalgebra-to-hopf-b.i002
		\item Secara dual, jika $A$ bialjabar komutatif di atas $\mathcal{V}$,
		maka $A\dcate{Comod}$ dan $\cated{Comod}A$ merupakan kategori monoidal
		simetris.
	\end{enumerate}
\end{proposition}

% segment-id: o014.aljabr2.chapter7.coalgebra-to-hopf-b.pf001
\begin{proof}
	Cukup bahas kasus kokomutatif. Dengan menggunakan notasi sebelumnya,
	pernyataan bahwa $c(M_1,M_2)$ merupakan morfisme dalam $A\dcate{Mod}$
	setara dengan membuktikan bahwa batas luar diagram berikut komutatif:
% segment-id: o014.aljabr2.chapter7.coalgebra-to-hopf-b.g001
	\[\begin{tikzcd}[row sep=large]
		A \otimes M_1 \otimes M_2 \arrow[r, "{\Delta \otimes \identity \otimes \identity}" inner sep=0.6em] \arrow[d, "{\identity \otimes c(M_1, M_2)}" description] & A \otimes A \otimes M_1 \otimes M_2 \arrow[r, "{\identity \otimes c(A, M_1) \otimes \identity}" inner sep=0.6em] \arrow[d, "{c(A, A) \otimes c(M_1, M_2)}" description] & A \otimes M_1 \otimes A \otimes M_2 \arrow[r, "{\mu_1 \otimes \mu_2}" inner sep=0.6em] \arrow[d, "{c(A \otimes M_1, A \otimes M_2)}" description] & M_1 \otimes M_2 \arrow[d, "{c(M_1, M_2)}" description] \\
		A \otimes M_2 \otimes M_1 \arrow[r, "{\Delta \otimes \identity \otimes \identity}"' inner sep=0.6em] & A \otimes A \otimes M_2 \otimes M_2 \arrow[r, "{\identity \otimes c(A, M_2) \otimes \identity}"' inner sep=0.6em] & A \otimes M_2 \otimes A \otimes M_1 \arrow[r, "{\mu_2 \otimes \mu_1}"' inner sep=0.6em] & M_2 \otimes M_1 .
	\end{tikzcd}\]

	Persegi di sebelah kiri komutatif karena $A$ kokomutatif, sedangkan persegi
	di sebelah kanan komutatif karena funktorialitas $c$. Dengan menggambar
	diagram anyamannya, terlihat bahwa
% segment-id: o014.aljabr2.chapter7.coalgebra-to-hopf-b.g002
	\begin{multline*}
		c(A \otimes M_1, A \otimes M_2) = \\
		(\identity_A \otimes c(A, M_2) \otimes \identity_{M_1}) (c(A, A) \otimes c(M_1, M_2)) (\identity_A \otimes c(A, M_1) \otimes \identity_{M_2}),
	\end{multline*}
	dan karena $\mathcal{V}$ kategori monoidal simetris, persegi di tengah juga
	komutatif. Semua sifat lain dari struktur monoidal simetris pada
	$A\dcate{Mod}$ cukup diverifikasi di $\mathcal{V}$.
\end{proof}

% segment-id: o014.aljabr2.chapter7.coalgebra-to-hopf-b.p002
Sekarang kita memperkenalkan konsep aljabar Hopf. Misalkan $\mathcal{V}$
kategori monoidal beranyam.

% segment-id: o014.aljabr2.chapter7.coalgebra-to-hopf-b.d001
\begin{definition}\label{def:Hopf-algebra}
% segment-id: o014.aljabr2.chapter7.coalgebra-to-hopf-b.x001
	\index{aljabar Hopf (Hopf algebra)}
% segment-id: o014.aljabr2.chapter7.coalgebra-to-hopf-b.x002
	\index{antipoda (antipode)}
	Misalkan data $(A,\mu,\eta,\Delta,\epsilon)$ merupakan bialjabar $A$ di
	atas $\mathcal{V}$. Jika suatu isomorfisme $S:A\rightiso A$ dalam
	$\mathcal{V}$ membuat diagram berikut komutatif,
% segment-id: o014.aljabr2.chapter7.coalgebra-to-hopf-b.g003
	\[\begin{tikzcd}
		A \otimes A \arrow[d, "\mu"'] & A \otimes A \arrow[r, "{S \otimes \identity_A}"] \arrow[l, "{\identity_A \otimes S}"'] & A \otimes A \arrow[d, "\mu"] \\
		A & A \arrow[u, "\Delta"] \arrow[r, "\eta\epsilon"'] \arrow[l, "\eta\epsilon"] & A ,
	\end{tikzcd}\]
	yaitu,
% segment-id: o014.aljabr2.chapter7.coalgebra-to-hopf-b.q001
	\[ \mu (\identity_A \otimes S) \Delta = \eta\epsilon = \mu (S \otimes \identity_A) \Delta, \]
	maka $S$ disebut \emph{antipoda} dari $A$. Bialjabar yang dilengkapi
	antipoda disebut \emph{aljabar Hopf}. Morfisme antaraljabar Hopf
	didefinisikan sebagai morfisme antarbialjabar yang
	mendasarinya.
\end{definition}

% segment-id: o014.aljabr2.chapter7.coalgebra-to-hopf-b.p003
Sebagai contoh, $\munit$ merupakan aljabar Hopf dengan antipoda
$\identity_{\munit}$.

% segment-id: o014.aljabr2.chapter7.coalgebra-to-hopf-b.p004
Jika ditafsirkan menggunakan konvolusi dalam
Catatan~\ref{rem:coalgebra-convolution}, antipoda $S$ tepat merupakan invers
dua sisi dari $\identity_A\in\End(A)$ terhadap $\star$. Antipoda dalam definisi
aljabar Hopf, jika ada, bersifat unik. Hal ini langsung diperoleh dengan
menerapkan hasil berikut pada $f=\identity_A$; hasil tersebut juga menunjukkan
bahwa morfisme selalu mempertahankan antipoda.

% segment-id: o014.aljabr2.chapter7.coalgebra-to-hopf-b.pr002
\begin{proposition}\label{prop:antipode-homo}
	Misalkan $A$ dan $A'$ aljabar Hopf di atas $\mathcal{V}$ dengan antipoda
	masing-masing $S$ dan $S'$. Untuk setiap morfisme aljabar Hopf
	$f:A\to A'$, berlaku $S'f=fS$.
\end{proposition}

% segment-id: o014.aljabr2.chapter7.coalgebra-to-hopf-b.pf002
\begin{proof}
	Berdasarkan funktorialitas konvolusi $\star$ dalam
	\eqref{eqn:convolution-functorial}, kedua pemetaan
% segment-id: o014.aljabr2.chapter7.coalgebra-to-hopf-b.g004
	\[\begin{tikzcd}
		\Hom(A', A') \arrow[r, "{g \mapsto gf}"] & \Hom(A, A')
	\end{tikzcd}, \quad \begin{tikzcd}
		\Hom(A, A) \arrow[r, "{h \mapsto fh}"] & \Hom(A, A')
	\end{tikzcd}\]
	merupakan homomorfisme monoid konvolusi. Jadi, $S'f$ dan $fS$ keduanya
	merupakan invers konvolusi dari $f\in\Hom(A,A')$.
\end{proof}

% segment-id: o014.aljabr2.chapter7.coalgebra-to-hopf-b.p005
Karena struktur anyaman pada $\mathcal{V}$ memberikan struktur anyaman pada
$\mathcal{V}^{\opp}$, konsep bialjabar jelas bersifat swadual. Hasil berikut
menunjukkan bahwa definisi antipoda juga demikian.

% segment-id: o014.aljabr2.chapter7.coalgebra-to-hopf-b.pr003
\begin{proposition}\label{prop:bialgebra-duality}
	Jika $(A,\mu,\eta,\Delta,\epsilon)$ bialjabar di atas $\mathcal{V}$, maka
	$(A,\eta,\mu,\epsilon,\Delta)$ merupakan bialjabar di atas
	$\mathcal{V}^{\opp}$. Jika aljabar Hopf
	$(A,\mu,\eta,\Delta,\epsilon)$ memiliki antipoda $S$, maka $S^{-1}$
	merupakan antipoda dari $(A,\eta,\mu,\epsilon,\Delta)$ sebagai bialjabar
	di atas $\mathcal{V}^{\opp}$; khususnya, objek terakhir ini merupakan
	aljabar Hopf di atas $\mathcal{V}^{\opp}$.
\end{proposition}

% segment-id: o014.aljabr2.chapter7.coalgebra-to-hopf-b.pf003
\begin{proof}
	Diagram dalam Definisi~\ref{def:Hopf-algebra} bersifat swadual, tetapi arah
	$S$ dibalik.
\end{proof}

% segment-id: o014.aljabr2.chapter7.coalgebra-to-hopf-b.pr004
\begin{proposition}\label{prop:Hopf-antiautomorphism}
% segment-id: o014.aljabr2.chapter7.coalgebra-to-hopf-b.x003
	\index[sym1]{Acop@$A_{\mathrm{cop}}, A^{\mathrm{op}}_{\mathrm{cop}}$}
	Untuk bialjabar $A$, sesuai dengan
	Definisi~\ref{def:opposite-braided-algebra},
% segment-id: o014.aljabr2.chapter7.coalgebra-to-hopf-b.l002
	\begin{compactitem}
% segment-id: o014.aljabr2.chapter7.coalgebra-to-hopf-b.i003
		\item membalik perkalian $A$ menggunakan struktur anyaman $c(A,A)$
		menghasilkan aljabar $A^{\mathrm{op}}$;
% segment-id: o014.aljabr2.chapter7.coalgebra-to-hopf-b.i004
		\item membalik komultiplikasi $A$ menggunakan $c(A,A)^{-1}$
		menghasilkan koaljabar $A_{\mathrm{cop}}$.
	\end{compactitem}
	Kedua operasi ini bersama-sama menghasilkan bialjabar
	$A^{\mathrm{op}}_{\mathrm{cop}}$. Jika $A$ memiliki antipoda $S$, maka $S$
	juga merupakan antipoda dari $A^{\mathrm{op}}_{\mathrm{cop}}$, dan dalam
	keadaan ini $S:A\rightiso A^{\mathrm{op}}_{\mathrm{cop}}$ merupakan
	isomorfisme bialjabar.
\end{proposition}

% segment-id: o014.aljabr2.chapter7.coalgebra-to-hopf-b.pf004
\begin{proof}
	Verifikasi rutin menunjukkan bahwa
	$A^{\mathrm{op}}_{\mathrm{cop}}$ memang bialjabar dengan antipoda $S$;
	perinciannya dihilangkan. Pernyataan yang tersisa terdiri atas dua bagian.
	Pertama, $S$ mempertahankan unit dan kounit. Ini mudah: dalam
	Proposisi~\ref{prop:antipode-homo}, ambil berturut-turut $A=\munit$ dan
	$A'=\munit$.

	Kedua, $S$ mempertahankan perkalian dan komultiplikasi. Bukti versi umumnya
	lebih rumit; lihat \cite[\S 9, Proposisi 2]{JS91} atau
	\cite[Proposisi 1.22]{AM10}. Di sini hanya diberikan garis besar kasus
	khusus $\mathcal{V}=\Bbbk\dcate{Mod}$,
	$\otimes=\otimes_{\Bbbk}$ dengan struktur anyaman standarnya. Inilah
	satu-satunya kasus yang akan diperlukan kemudian.

	Ingat dari \S\ref{sec:alg-in-monoidal-cat} bahwa $A\otimes A$ secara alami
	merupakan koaljabar. Lengkapi $\Hom(A\otimes A,A)$ dengan konvolusi dari
	Catatan~\ref{rem:coalgebra-convolution}, yang ditulis $\ast$, lalu tinjau
	unsur-unsurnya
% segment-id: o014.aljabr2.chapter7.coalgebra-to-hopf-b.q002
	\[ f := \mu, \quad g := \mu c(A, A) (S \otimes S), \quad h := S \mu. 
	\]
	Cukup dibuktikan bahwa
	$h\ast f=\eta_A\epsilon_{A\otimes A}=f\ast g$. Kesamaan ini akan
	menunjukkan bahwa $f$ merupakan unsur invertibel dalam monoid
	$(\Hom(A\otimes A,A),\ast)$, dan karena itu $h=g$. Misalkan $a,b\in A$,
	dan tuliskan ekspansi
% segment-id: o014.aljabr2.chapter7.coalgebra-to-hopf-b.g005
	\begin{gather*}
		\Delta(a) = \sum_i a_i^{(1)} \otimes a_i^{(2)}, \quad
		\Delta(b) = \sum_j b_j^{(1)} \otimes b_j^{(2)}.
	\end{gather*}
	Tuliskan $\mu$ langsung sebagai perkalian. Berdasarkan struktur anyaman
	standar pada $\Bbbk\dcate{Mod}$, diperoleh
% segment-id: o014.aljabr2.chapter7.coalgebra-to-hopf-b.g006
	\begin{align*}
		(h \ast f)(a \otimes b) & = \sum_{i, j} h\left(a^{(1)}_i \otimes b^{(1)}_j\right) f\left(a^{(2)}_i \otimes b^{(2)}_j\right) \\
		& = \sum_{i, j} S\left( a^{(1)}_i b^{(1)}_j \right) a^{(2)}_i b^{(2)}_j \\
		& = (\underbracket{S \star \identity_A}_{\text{konvolusi pada }\End(A)})(ab) = \eta_A \epsilon_A (ab), \\
		(f \ast g)(a \otimes b) & = \cdots = \sum_{i, j} a^{(1)}_i b^{(1)}_j S\left( b^{(2)}_j \right) S\left( a^{(2)}_i \right) \\
		& = \sum_i a^{(1)}_i \eta_A \epsilon_A(b) S\left( a^{(2)}_i \right) \\
		& = \eta_A \epsilon_A (a) \cdot \eta_A \epsilon_A (b) = \eta_A \epsilon_A(ab), \\
		\eta_A \epsilon_{A \otimes A}(a \otimes b) & = \eta_A\left( \epsilon_A(a)\epsilon_A(b) \right) = \eta_A \epsilon_A(ab),
	\end{align*}
	seperti yang diklaim. Jadi, $S$ mempertahankan perkalian bialjabar. Argumen
	serupa menunjukkan bahwa $S$ mempertahankan komultiplikasi.
\end{proof}

% segment-id: o014.aljabr2.chapter7.coalgebra-to-hopf-b.p006
Khususnya, jika $\mathcal{V}$ kategori monoidal simetris, maka $S^2$ merupakan
automorfisme dari $A$.

% segment-id: o014.aljabr2.chapter7.coalgebra-to-hopf-b.c001
\begin{corollary}\label{prop:Hopf-involution}
	Misalkan aljabar Hopf $A$ memiliki antipoda $S$. Jika $A$ komutatif atau
	kokomutatif, maka $S^2=\identity$.
\end{corollary}

% segment-id: o014.aljabr2.chapter7.coalgebra-to-hopf-b.pf005
\begin{proof}
	Terhadap konvolusi $\star$ pada $\End(A)$
	(Catatan~\ref{rem:coalgebra-convolution}), kita mempunyai
% segment-id: o014.aljabr2.chapter7.coalgebra-to-hopf-b.q003
	\[ S \star S^2 = \mu (S \otimes S^2) \Delta = \mu (S \otimes S)(\identity \otimes S) \Delta. \]

	Jika $A$ komutatif, maka $\mu c(A,A)=\mu$, dan
	Proposisi~\ref{prop:Hopf-antiautomorphism} mengubah ruas di atas menjadi
% segment-id: o014.aljabr2.chapter7.coalgebra-to-hopf-b.q004
	\[ \mu c(A, A) (S \otimes S) (\identity \otimes S)\Delta = S\mu (\identity \otimes S)\Delta = S(\identity \star S). \]
	Definisi antipoda menunjukkan bahwa
	$\identity\star S=\eta\epsilon$, sedangkan proposisi tersebut memberikan
	$S\eta=\eta$; hasil akhirnya ialah $\eta\epsilon$.

	Jika $A$ kokomutatif, maka $c(A,A)\Delta=\Delta$, dan funktorialitas struktur
	anyaman mengubah $S\star S^2$ menjadi
% segment-id: o014.aljabr2.chapter7.coalgebra-to-hopf-b.q005
	\[ \mu (S \otimes S)(\identity \otimes S) c(A, A) \Delta = \mu c(A, A) (S \otimes S)(S \otimes \identity)\Delta, \]
	yang oleh argumen yang sama kemudian berubah menjadi
	$S(S\star\identity)=\eta\epsilon$.

	Jadi, $S^2$ merupakan invers konvolusi dari $S$, sehingga
	$S^2=\identity$.
\end{proof}

% segment-id: o014.aljabr2.chapter7.coalgebra-to-hopf-b.p007
Semua contoh berikut mengambil $\mathcal{V}=\Bbbk\dcate{Mod}$ dan
$\otimes=\otimes_{\Bbbk}$, dengan $\Bbbk$ gelanggang komutatif.

% segment-id: o014.aljabr2.chapter7.coalgebra-to-hopf-b.ex001
\begin{example}\label{eg:group-Hopf-algebra}
	Misalkan $G$ grup. Bentuk aljabar grup-$\Bbbk$ $\Bbbk[G]$.
	Contoh~\ref{eg:monoid-bialgebra} memberi $\Bbbk[G]$ struktur bialjabar.
	Definisikan automorfisme modul-$\Bbbk$
	$S:\Bbbk[G]\to\Bbbk[G]$ melalui
% segment-id: o014.aljabr2.chapter7.coalgebra-to-hopf-b.q006
	\[ S(g) = g^{-1}, \quad g \in G. \]
	Mudah diverifikasi bahwa $S$ memenuhi diagram komutatif yang disyaratkan
	bagi antipoda. Jadi, $\Bbbk[G]$ merupakan aljabar Hopf.
\end{example}

% segment-id: o014.aljabr2.chapter7.coalgebra-to-hopf-b.ex002
\begin{example}[H.\ Hopf]
% segment-id: o014.aljabr2.chapter7.coalgebra-to-hopf-b.x004
	\index{ruang-H, grup-H (H-space, H-group)}
	Tuliskan $\{\mathrm{pt}\}$ untuk himpunan satu titik. Misalkan ruang
	topologis $X$ dilengkapi pemetaan kontinu
	$m:X\times X\to X$ (perkalian) dan
	$e:\{\mathrm{pt}\}\to X$ (setara dengan memilih suatu unsur dari $X$,
	yaitu unit, yang juga ditulis $e$), sedemikian sehingga kedua pemetaan
	memenuhi sifat-sifat yang diperlukan bagi asosiativitas perkalian dan unit,
	hingga homotopi. Maka $(X,m,e)$ disebut \emph{ruang-H}. Jika juga dipilih
	pemetaan kontinu $i:X\to X$ (pengambilan invers) yang, hingga homotopi,
	memenuhi sifat-sifat yang diperlukan bagi invers, maka $(X,m,e,i)$ disebut
	\emph{grup-H}.

	Monoid topologis (atau grup topologis) tentu merupakan ruang-H (atau
	grup-H), tetapi syarat ruang-H atau grup-H jauh lebih longgar. Contoh
	alaminya ialah ruang gelung. Untuk ruang topologis $M$ dengan titik pangkal
	terpilih $x_0\in M$, pemetaan kontinu $\gamma:[0,1]\to M$ yang memenuhi
	$\gamma(0)=x_0=\gamma(1)$ disebut \emph{gelung} di $M$. Semua gelung secara
	alami membentuk ruang topologis $\Omega(M,x_0)$. Jika $m$ diambil sebagai
	penyambungan ujung ke pangkal dua gelung, $e$ sebagai pemetaan konstan
	$x_0$, dan $i$ sebagai pembalikan arah gelung, maka secara intuitif
	$\Omega(M,x_0)$ membentuk grup-H. Kuncinya ialah asosiativitas dan sifat
	invers hanya berlaku dalam arti homotopi.

	Selanjutnya, tinjau subkategori penuh $\cate{S}$ dari $\cate{Top}$ yang
	terdiri atas semua kompleks CW\footnote{Ini adalah istilah topologis, yang
	agak berbeda dari kompleks yang didefinisikan dalam aljabar linear.};
	kategori ini memuat ruang-ruang yang umum dijumpai dalam geometri, seperti
	manifold. Misalkan $\Bbbk$ medan. Funktor kohomologi berkoefisien
	$\Bbbk$, yakni $\Hm^\bullet$, merupakan funktor dari $\cate{S}^{\opp}$ ke
	kategori ruang vektor-$\Bbbk$ bergradasi
	$\cate{Vect}(\Bbbk)^{\Z_{\geq0}}$. Lengkapi
	$\cate{Vect}(\Bbbk)^{\Z_{\geq0}}$ dengan struktur anyaman yang ditentukan
	oleh aturan tanda Koszul \eqref{eqn:Koszul-braiding} (dalam notasi di sana,
	ambil $\epsilon(a)=a\bmod\;2$). Hasil kali cawan
	$\mu:=\cup$ pada kohomologi membuat funktor $\Hm^\bullet$ berfaktor melalui
	$\cate{CAlg}\left(\cate{Vect}(\Bbbk)^{\Z_{\geq0}}\right)$; singkatnya,
	kohomologi merupakan aljabar komutatif di atas
	$\cate{Vect}(\Bbbk)^{\Z_{\geq0}}$.

	Di sisi lain, beri $\cate{S}$ struktur monoidal yang ditentukan oleh produk
	ruang, dengan $\{\mathrm{pt}\}$ sebagai unit. Rumus Künneth dalam topologi
	menunjukkan bahwa $\Hm^\bullet$ merupakan funktor monoidal. Selain itu,
	pemetaan-pemetaan yang homotopik menginduksi pemetaan yang sama pada tingkat
	kohomologi. Oleh karena itu, jika $X\in\Obj(\cate{S})$ merupakan ruang-H,
	maka $\Hm^\bullet(X)$ juga merupakan koaljabar di atas
	$\cate{Vect}(\Bbbk)^{\Z_{\geq0}}$.

	Pembaca yang mengenal topologi aljabar akan mudah menyimpulkan bahwa
	struktur aljabar dan koaljabar pada $\Hm^\bullet(X)$ kompatibel, sehingga
	memberikan bialjabar di atas
	$\cate{Vect}(\Bbbk)^{\Z_{\geq0}}$. Jika lebih lanjut $X$ merupakan grup-H
	dan automorfisme yang diinduksi oleh pemetaan $i$ ditulis
	$S\in\Aut(\Hm^\bullet(X))$, maka $\Hm^\bullet(X)$ menjadi aljabar Hopf.

	Sesungguhnya, perkalian (hasil kali cawan) pada $\Hm^\bullet(X)$ tidak lain
	adalah citra penyertaan diagonal
	$\mathrm{diag}:X\to X\times X$ melalui rumus Künneth, sedangkan unitnya
	adalah citra $X\to\{\mathrm{pt}\}$. Karena itu, semua sifat yang
	diperlukan bagi bialjabar dan antipoda dapat diverifikasi pada tingkat ruang
	topologis, hingga homotopi. Sebagai contoh, identitas antipoda
	$\mu(S\otimes\identity)\Delta=\eta\epsilon$ berasal dari
% segment-id: o014.aljabr2.chapter7.coalgebra-to-hopf-b.q007
	\[ \left[ X \xrightarrow{\mathrm{diag}} X \times X \xrightarrow{(i, \identity)} X \times X \xrightarrow{m} X\right] \;\text{ekuivalen secara homotopi dengan}\; \left[ X \to \{\mathrm{pt}\} \xrightarrow{e} X \right]. \]

	Dengan demikian, kohomologi suatu grup-H secara alami merupakan aljabar Hopf
	komutatif dengan struktur yang kaya. Inilah motivasi awal Hopf untuk
	mempelajari aljabar semacam itu.
\end{example}

% segment-id: o014.aljabr2.chapter7.coalgebra-to-hopf-b.p008
Latihan-latihan akan memberikan lebih banyak contoh aljabar Hopf. Kebanyakan
contoh itu komutatif atau kokomutatif, tetapi terdapat pula banyak aljabar Hopf
yang tidak memiliki kedua sifat tersebut. Kelas terpenting di antaranya ialah
\emph{grup kuantum}. Pembahasan terkait memerlukan landasan teori Lie dan tidak
dibahas dalam buku ini.
